\section{Realization of Root Systems}

Let
\[
    E=\mathbb{R}^n
\]
with the usual inner product, and let
\[
    w_1,\ldots,w_n
\]
be the usual orthonormal basis of \(\mathbb{R}^n\). The \(\mathbb{Z}\)-span
of this basis is a lattice, denoted by
\[
    I=\sum_{i=1}^n \mathbb{Z}w_i .
\]
The root systems below will be defined as sets of vectors in a suitable lattice
having prescribed length or lengths.

\begin{remark}
    To check that such a set \(\Delta\) is a root system, one usually verifies the
    following points. First, \(\Delta\) is finite, spans the ambient Euclidean space,
    and does not contain \(0\). Also,
    \[
        \mathbb{R}\alpha\cap \Delta=\{\alpha,-\alpha\}
        \qquad
        \text{for all } \alpha\in\Delta .
    \]
    Moreover, for \(\alpha,\beta\in\Delta\),
    \[
        \langle \beta,\alpha^\vee\rangle
        =
        \frac{2(\beta,\alpha)}{(\alpha,\alpha)}
        \in \mathbb{Z}.
    \]
    Finally, one checks that the reflection
    \[
        s_\alpha(\beta)
        =
        \beta-\frac{2(\beta,\alpha)}{(\alpha,\alpha)}\alpha
    \]
    maps \(\Delta\) into itself. Since \(s_\alpha\) preserves lengths, it is often
    enough to check that \(s_\alpha(\beta)\) lies in the relevant lattice.
\end{remark}

\subsection{\texorpdfstring{The root system $A_n$}{The root system An}}

Let \(E\) be the \(n\)-dimensional subspace of \(\mathbb{R}^{n+1}\) orthogonal
to
\[
    w_1+\cdots+w_{n+1}.
\]
Set
\[
    I'=I\cap E
    =
    \left\{
    \sum_{i=1}^{n+1} k_iw_i
    \ \middle|\
    k_i\in\mathbb{Z},\ \sum_{i=1}^{n+1}k_i=0
    \right\}.
\]
Define
\[
    A_n
    =
    \left\{
    \alpha\in I'
    \ \middle|\
    (\alpha,\alpha)=2
    \right\}.
\]
Then
\[
    A_n
    =
    \{\,w_i-w_j\mid 1\le i\ne j\le n+1\,\}.
\]

A choice of simple roots is
\[
    \alpha_1=w_1-w_2,\quad
    \alpha_2=w_2-w_3,\quad
    \ldots,\quad
    \alpha_n=w_n-w_{n+1}.
\]
Thus the Dynkin diagram is
\[
    \begin{tikzpicture}[rootdiagram]
        \node[vertex,label=below:{\(\alpha_1=w_1-w_2\)}] (a1) at (0,0) {};
        \node[vertex,label=below:{\(\alpha_2=w_2-w_3\)}] (a2) at (2.2,0) {};
        \node at (3.5,0) {\(\cdots\)};
        \node[vertex,label=below:{\(\alpha_n=w_n-w_{n+1}\)}] (an) at (5.0,0) {};
        \draw (a1)--(a2);
        \draw (a2)--(3.0,0);
        \draw (4.0,0)--(an);
    \end{tikzpicture}
\]
The corresponding positive and negative roots are
\[
    \Delta^+
    =
    \{\,w_i-w_j\mid i<j\,\},
    \qquad
    \Delta^-
    =
    \{\,w_i-w_j\mid i>j\,\}.
\]

For \(\alpha_i=w_i-w_{i+1}\), the reflection \(s_{\alpha_i}\) interchanges the
subscripts \(i\) and \(i+1\), and leaves all other subscripts fixed. Hence
\(s_{\alpha_i}\) corresponds to the transposition \((i\ i+1)\) in \(S_{n+1}\).
Since the Weyl group is generated by the simple reflections, we obtain
\[
    W(A_n)\cong S_{n+1}.
\]

\subsection{\texorpdfstring{The root system \(D_n\)}{The root system Dn}}

Let \(E=\mathbb{R}^n\). Define
\[
    D_n
    =
    \left\{
    \alpha\in I
    \ \middle|\
    (\alpha,\alpha)=2
    \right\}.
\]
Then
\[
    D_n
    =
    \{\,\pm w_i\pm w_j\mid 1\le i<j\le n\,\}.
\]

A choice of simple roots is
\[
    \alpha_1=w_1-w_2,\quad
    \alpha_2=w_2-w_3,\quad
    \ldots,\quad
    \alpha_{n-1}=w_{n-1}-w_n,\quad
    \alpha_n=w_{n-1}+w_n.
\]
The Dynkin diagram is
\[
    \begin{tikzpicture}[rootdiagram]
        \node[vertex,label=below:{\(w_1-w_2\)}] (a1) at (0,0) {};
        \node[vertex,label=below:{\(w_2-w_3\)}] (a2) at (1.8,0) {};
        \node at (3.0,0) {\(\cdots\)};
        \node[vertex,label=below:{\(w_{n-2}-w_{n-1}\)}] (a3) at (4.3,0) {};
        \node[vertex,label=above right:{\(w_{n-1}-w_n\)}] (a4) at (5.8,0.7) {};
        \node[vertex,label=below right:{\(w_{n-1}+w_n\)}] (a5) at (5.8,-0.7) {};

        \draw (a1)--(a2);
        \draw (a2)--(2.55,0);
        \draw (3.45,0)--(a3);
        \draw (a3)--(a4);
        \draw (a3)--(a5);
    \end{tikzpicture}
\]

\subsection{\texorpdfstring{The root system \(B_n\)}{The root system Bn}}

Let \(E=\mathbb{R}^n\). Define
\[
    B_n
    =
    \left\{
    \alpha\in I
    \ \middle|\
    (\alpha,\alpha)=1 \text{ or } 2
    \right\}.
\]
Then
\[
    B_n
    =
    \{\,\pm w_i,\ \pm w_i\pm w_j
    \mid 1\le i<j\le n\,\}.
\]

A choice of simple roots is
\[
    \alpha_1=w_1-w_2,\quad
    \alpha_2=w_2-w_3,\quad
    \ldots,\quad
    \alpha_{n-1}=w_{n-1}-w_n,\quad
    \alpha_n=w_n.
\]
Notice that every short root is expressible in this basis. For example,
\[
    w_i
    =
    (w_i-w_{i+1})+(w_{i+1}-w_{i+2})
    +\cdots+(w_{n-1}-w_n)+w_n.
\]
Similarly, the long roots \(w_i-w_j\) and \(w_i+w_j\) are expressible in
this basis.

The Dynkin diagram is
\[
    \begin{tikzpicture}[rootdiagram]
        \node[vertex,label=below:{\(w_1-w_2\)}] (a1) at (0,0) {};
        \node[vertex,label=below:{\(w_2-w_3\)}] (a2) at (1.8,0) {};
        \node at (3.0,0) {\(\cdots\)};
        \node[vertex,label=below:{\(w_{n-1}-w_n\)}] (a3) at (4.4,0) {};
        \node[vertex,label=below:{\(w_n\)}] (a4) at (6.0,0) {};

        \draw (a1)--(a2);
        \draw (a2)--(2.55,0);
        \draw (3.45,0)--(a3);
        \draw[double, double distance=2pt, dynarrow] (a3)--(a4);
    \end{tikzpicture}
\]
The arrow points toward the shorter root \(w_n\).

\subsection{\texorpdfstring{The root system \(C_n\)}{The root system Cn}}

The root system \(C_n\) is dual to \(B_n\). Thus
\[
    C_n
    =
    \{\,\pm 2w_i,\ \pm w_i\pm w_j
    \mid 1\le i<j\le n\,\}.
\]
Equivalently, it is obtained from \(B_n\) by replacing each root
\(\alpha\) by its coroot
\[
    \alpha^\vee=\frac{2\alpha}{(\alpha,\alpha)}.
\]

A choice of simple roots is
\[
    \alpha_1=w_1-w_2,\quad
    \alpha_2=w_2-w_3,\quad
    \ldots,\quad
    \alpha_{n-1}=w_{n-1}-w_n,\quad
    \alpha_n=2w_n.
\]
The Dynkin diagram is
\[
    \begin{tikzpicture}[rootdiagram]
        \node[vertex,label=below:{\(w_1-w_2\)}] (a1) at (0,0) {};
        \node[vertex,label=below:{\(w_2-w_3\)}] (a2) at (1.8,0) {};
        \node at (3.0,0) {\(\cdots\)};
        \node[vertex,label=below:{\(w_{n-1}-w_n\)}] (a3) at (4.4,0) {};
        \node[vertex,label=below:{\(2w_n\)}] (a4) at (6.0,0) {};

        \draw (a1)--(a2);
        \draw (a2)--(2.55,0);
        \draw (3.45,0)--(a3);
        \draw[double, double distance=2pt, dynarrowleft] (a3)--(a4);
    \end{tikzpicture}
\]
Here the arrow points toward the shorter root \(w_{n-1}-w_n\).

\subsection{\texorpdfstring{The root system \(F_4\)}{The root system F4}}

Let \(E=\mathbb{R}^4\), and set
\[
    I'
    =
    I+\frac{\mathbb{Z}}{2}(w_1+w_2+w_3+w_4).
\]
Define
\[
    F_4
    =
    \left\{
    \alpha\in I'
    \ \middle|\
    (\alpha,\alpha)=1 \text{ or } 2
    \right\}.
\]
Then
\[
    F_4
    =
    \left\{
    \pm w_i,\
    \pm w_i\pm w_j,\
    \frac12(\pm w_1\pm w_2\pm w_3\pm w_4)
    \ \middle|\
    1\le i<j\le 4
    \right\}.
\]

A choice of simple roots is
\[
    \alpha_1=w_1-w_2,\qquad
    \alpha_2=w_2-w_3,\qquad
    \alpha_3=w_3,\qquad
    \alpha_4=-\frac12(w_1+w_2+w_3+w_4).
\]
The Dynkin diagram is
\[
    \begin{tikzpicture}[rootdiagram]
        \node[vertex,label=below:{\(w_1-w_2\)}] (a1) at (0,0) {};
        \node[vertex,label=below:{\(w_2-w_3\)}] (a2) at (1.8,0) {};
        \node[vertex,label=below:{\(w_3\)}] (a3) at (3.4,0) {};
        \node[vertex,label=below:{\(-\frac12(w_1+w_2+w_3+w_4)\)}] (a4) at (5.4,0) {};

        \draw (a1)--(a2);
        \draw[double, double distance=2pt, dynarrow] (a2)--(a3);
        \draw (a3)--(a4);
    \end{tikzpicture}
\]

\subsection{\texorpdfstring{The root system \(G_2\)}{The root system G2}}

There is also an exceptional realization of \(G_2\) in terms of the octonions:
\[
    G_2=\operatorname{Der}(\mathbb{O}).
\]
Over \(\mathbb{C}\), one obtains the complex Lie algebra of type \(G_2\) by
complexifying:
\[
    \operatorname{Der}(\mathbb{O}\otimes_{\mathbb{R}}\mathbb{C}).
\]

For the abstract root system, take \(E\) to be the \(2\)-dimensional subspace
of \(\mathbb{R}^3\) orthogonal to
\[
    w_1+w_2+w_3.
\]
Let
\[
    I'
    =
    I\cap E.
\]
Define
\[
    G_2
    =
    \left\{
    \alpha\in I'
    \ \middle|\
    (\alpha,\alpha)=2 \text{ or } 6
    \right\}.
\]
Then
\[
    G_2
    =
    \{\,\pm(w_i-w_j)\mid 1\le i<j\le 3\,\}
    \cup
    \{\,\pm(2w_i-w_j-w_k)
    \mid \{i,j,k\}=\{1,2,3\}\,\}.
\]

A choice of simple roots is
\[
    \alpha_1=w_1-w_2,
    \qquad
    \alpha_2=-2w_1+w_2+w_3.
\]
The Dynkin diagram is
\[
    \begin{tikzpicture}[rootdiagram]
        \node[vertex,label=below:{\(w_1-w_2\)}] (a1) at (0,0) {};
        \node[vertex,label=below:{\(-2w_1+w_2+w_3\)}] (a2) at (3.0,0) {};

        \draw (a1)--(a2);
        \draw ([yshift=2pt]a1)--([yshift=2pt]a2);
        \draw[dynarrowleft] ([yshift=-2pt]a1)--([yshift=-2pt]a2);
    \end{tikzpicture}
\]
The arrow points toward the shorter root \(w_1-w_2\).




